% Topologia de despliegue: un emisor, muchos receptores, todo el calculo en el servidor.
%
% Se dibuja en TikZ y no se importa como imagen para que use la tipografia del documento y
% escale sin perdida. Pensada para leerse tambien en escala de grises: las cajas se
% distinguen por trazo y relleno claro, nunca por color.
\begin{figure}[!ht]
\centering
\begin{tikzpicture}[
    font=\footnotesize,
    caja/.style     = {draw, rounded corners=2pt, align=center, fill=white,
                       minimum height=7mm, inner xsep=5pt},
    equipo/.style   = {caja, fill=black!4, minimum width=27mm},
    servidor/.style = {caja, fill=black!8, minimum width=40mm, font=\footnotesize\bfseries},
    zona/.style     = {draw, dashed, rounded corners=4pt, inner sep=7pt},
    flecha/.style   = {-{Stealth[length=2mm]}, thick},
    nota/.style     = {align=left, font=\scriptsize\itshape, text=black!65},
]

% ---------------- puesto del docente ----------------
\node[equipo] (mic)  {Micrófono de solapa};
\node[equipo, below=5mm of mic] (nav) {Navegador \texttt{/broadcast}\\[1pt]
                                       \scriptsize captura y trocea};
\draw[flecha] (mic) -- (nav);
\node[zona, fit=(mic)(nav),
      label={[font=\scriptsize\bfseries, yshift=-1pt]above:PUESTO DEL DOCENTE}]
      (zdoc) {};

% ---------------- servidor del aula ----------------
\node[servidor, right=22mm of nav.east, anchor=west] (app) {Aplicación .NET};
\node[servidor, below=7mm of app] (asr) {Servicio ASR (Python)};
\node[caja, below=4mm of asr, minimum width=40mm, fill=white] (gpu)
      {Tarjeta gráfica\\[1pt]\scriptsize modelo cargado una sola vez};
\draw[flecha] (asr) -- (gpu);
\draw[flecha] ([xshift=-8mm]app.south) -- ([xshift=-8mm]asr.north)
      node[midway, left, font=\scriptsize] {segmento};
\draw[flecha] ([xshift=8mm]asr.north) -- ([xshift=8mm]app.south)
      node[midway, right, font=\scriptsize] {texto};
\node[zona, fit=(app)(asr)(gpu),
      label={[font=\scriptsize\bfseries]above:SERVIDOR DEL AULA}] (zsrv) {};

% ---------------- alumnado ----------------
\node[equipo, right=22mm of app.east, anchor=west, yshift=6mm] (al1)
      {Navegador \texttt{/view}};
\node[equipo, below=3mm of al1] (al2) {Navegador \texttt{/view}};
\node[below=2mm of al2, font=\scriptsize] (aln) {$\cdots$ \emph{n} dispositivos};
\node[zona, fit=(al1)(al2)(aln),
      label={[font=\scriptsize\bfseries]above:ALUMNADO}] (zal) {};

% ---------------- flujo principal ----------------
\draw[flecha] (nav.east) -- node[above, font=\scriptsize] {audio}
                            node[below, font=\scriptsize] {16 kHz} (app.west);
\draw[flecha] (app.east) -- node[above, pos=0.22, font=\scriptsize] {subtítulos} (al1.west);
\draw[flecha] (app.east) -- (al2.west);

% ---------------- frontera de red ----------------
% El relleno superior es mayor que el resto: las etiquetas de cada zona se dibujan FUERA
% de su nodo, y con un relleno uniforme quedarian pisando el borde de esta caja.
\begin{scope}[on background layer]
  \node[draw, rounded corners=5pt, inner xsep=13pt, inner ysep=13pt,
        fill=black!2, fit=(zdoc)(zsrv)(zal)] (redint) {};
  \node[draw, rounded corners=5pt, inner xsep=0pt, inner ysep=0pt, fill=none,
        fit=(redint)] (red) {};
\end{scope}
% La etiqueta va abajo: arriba compite con los rotulos de las tres zonas.
\node[anchor=south west, font=\scriptsize\bfseries, text=black!70]
      at ([xshift=4pt, yshift=3pt]red.south west) {RED LOCAL DEL CENTRO};

% ---------------- lo que queda fuera ----------------
\node[nota, below=3mm of red.south, anchor=north] (fuera)
      {Toda conexión ajena a un rango privado se rechaza. El audio no sale del centro\\
       y no se almacena: la transcripción vive en memoria mientras dura la clase.};

\end{tikzpicture}
\caption{Topología del sistema desplegado. El reconocimiento ocurre \textbf{una sola vez en
el servidor} con independencia del número de dispositivos conectados; el alumnado solo
visualiza, sin instalar nada ni conceder permisos.}
\label{fig:topologia}
\end{figure}
